\documentclass[10pt]{article}
\usepackage[T1]{fontenc}
\usepackage[utf8]{inputenc}
\usepackage[french]{babel}
\usepackage{amsmath,amssymb,mathrsfs,mathpazo}
\usepackage{geometry}
\geometry{a4paper,margin=2cm}
\usepackage{graphicx,array,tabularx,tabularray,colortbl,multirow,multicol,booktabs,enumitem,pifont,ulem,fancybox,diagbox,cancel}
\usepackage{pstricks,pstricks-add,pst-plot,pst-tree,pst-node,pst-text,pst-func,pst-math,pst-bspline,pst-3dplot,pst-eucl,pst-solides3d}
\let\pscalculate\undefined
\usepackage{pst-bezier}
\usepackage{tikz,pgfplots,tkz-tab}
\usetikzlibrary{arrows,arrows.meta,patterns,calc,positioning,trees,decorations.markings,shapes.geometric,shapes.misc,shapes.arrows,intersections,decorations.pathmorphing,decorations.pathreplacing,backgrounds,shadows,babel}
\pgfplotsset{compat=1.18}
\usepgfplotslibrary{fillbetween}
\usepackage[np]{numprint}
\usepackage{listings,makecell,pas-tableur,twoopt}
\let\barre\undefined
\usepackage[pstricks]{tablvar}
\definecolor{darkgreen}{RGB}{0,100,0}
\providecommand{\red}{\color{red}}
\providecommand{\blue}{\color{blue}}
\providecommand{\green}{\color{green}}
\providecommand{\black}{\color{black}}
\providecommand{\white}{\color{white}}
\providecommand{\gray}{\color{gray}}
\providecommand{\cyan}{\color{cyan}}
\providecommand{\magenta}{\color{magenta}}
\providecommand{\yellow}{\color{yellow}}
\providecommand{\violet}{\color{violet}}
\providecommand{\orange}{\color{orange}}
\providecommand{\euro}{\texteuro}
\providecommand{\EUR}{\texteuro}
\providecommand{\ieme}{\textsuperscript{e}}
\providecommand{\degres}{\ensuremath{^\circ}}
\providecommand{\degree}{\ensuremath{^\circ}}
\providecommand{\vv}[1]{\overrightarrow{#1}}
\providecommand{\Coord}[1]{#1}
\providecommand{\decosix}{}
\providecommand{\starredbullet}{\ensuremath{\bullet}}
\pagestyle{empty}
\setlength{\parindent}{0pt}
\usepackage{adjustbox}\DeclareUnicodeCharacter{2236}{:}\DeclareUnicodeCharacter{2190}{\ensuremath{\leftarrow}}
\providecommand{\R}{\mathbb{R}}
\providecommand{\N}{\mathbb{N}}
\providecommand{\Z}{\mathbb{Z}}
\providecommand{\Q}{\mathbb{Q}}
\providecommand{\D}{\mathbb{D}}
\providecommand{\C}{\mathbb{C}}
\providecommand{\e}{{\mathrm{e}}}
\providecommand{\vect}[1]{\overrightarrow{#1}}
\providecommand{\vectt}[1]{\overrightarrow{\mathrm{#1}}}
\providecommand{\vvt}[1]{\overrightarrow{\mathrm{#1}}}
\providecommand{\np}[2][]{{#2~#1}}
\providecommand{\vv}[1]{{\overrightarrow{#1}}}
\renewcommand{\d}{{\mathrm{d}}}
\providecommand{\degres}{{^\circ}}
\providecommand{\degre}{{^\circ}}
\providecommand{\blue}{\color{blue}}
\providecommand{\red}{\color{red}}
\providecommand{\enskip}{\quad}
\providecommand{\llbracket}{\lbrack\!\lbrack}
\providecommand{\rrbracket}{\rbrack\!\rbrack}
\providecommand{\euro}{{\text{€}}}
\providecommand{\no}{{\text{n°}}}
\providecommand{\interval}[2]{{[#1;#2]}}
\providecommand{\intervFO}[2]{{[#1;#2[}}
\providecommand{\dotfill}{\ldots}

\providecommand{\ds}{\displaystyle}
\providecommand{\textding}[1]{#1}
\providecommand{\Oij}{\(\mathrm{O};\vec{\imath},\vec{\jmath}\)}
\providecommand{\Oijk}{\(\mathrm{O};\vec{\imath},\vec{\jmath},\vec{k}\)}
\providecommand{\Ouv}{\(\mathrm{O};\vec{u},\vec{v}\)}
\providecommand{\intervFF}[2]{\ensuremath{[#1;#2]}}
\renewcommand{\euro}{\ifmmode\text{\texteuro}\else\texteuro\fi}
\DeclareUnicodeCharacter{00B0}{\ifmmode^{\circ}\else\textdegree\fi}
\DeclareUnicodeCharacter{20AC}{\ifmmode\text{\texteuro}\else\texteuro\fi}
% Énoncés : Éducation nationale. Transcription et corrigés : APMEP, Denis Vergès et contributeurs. Découpage et classement : CoopMaths, Sébastien Lozano, Éric Elter et contributeurs. Adaptation de la présentation : Maths Explorer.
% CC BY-SA : https://coopmaths.fr/www/about/
% https://coopmaths.fr/alea/static/bac/2022/tex/bac_2022_05_sujet1_asie_4_cor.tex
\begin{document}
\noindent\textit{APMEP / CoopMaths. CC BY-SA. Présentation adaptée par Maths Explorer.}\par\medskip
\medskip

\emph{Principaux domaines abordés}: Étude de fonctions. Fonction logarithme.

\bigskip

Soit $f$ une fonction définie et dérivable sur $\R$. On considère les points A(1~;~3) et B(3~;~5).

On donne ci-dessous $\mathcal{C}_f$ la courbe représentative de $f$ dans un repère orthogonal du plan, ainsi que la tangente (AB) à la courbe $\mathcal{C}_f$ au point A.

\medskip

\begin{center}
\psset{unit=0.8cm,arrowsize=2pt 3}
\begin{pspicture*}(-7,-2.25)(8,6.25)
\psgrid[gridlabels=0pt,subgriddiv=1,gridwidth=0.15pt]
\psaxes[linewidth=1.25pt,labelFontSize=\scriptstyle]{->}(0,0)(-7,-2.25)(8,6.25)
\psplot[plotpoints=2000,linewidth=1.25pt,linecolor=red]{-7}{8}{x dup mul 1 add ln 3 add 2 ln sub}
\psplot[linestyle=dashed]{-7}{8}{x 2 add}
\uput[dr](1,3){A}\uput[ul](3,5){B}
\psdots[dotstyle=+,dotangle=45,dotscale=1.5](1,3)(3,5)
\uput[d](-5.5,5.75){\red $\mathcal{C}_f$}
\end{pspicture*}
\end{center}

\emph{Les trois parties de l'exercice peuvent être traitées de manière indépendante.}

\bigskip

\textbf{Partie A}

\medskip

\begin{enumerate}
\item 
On lit sur le graphique : $f(1) = 3$ et $f'(1) = 1$ (nombre dérivé égal au coefficient directeur de la droite (AB)).
\item 
	\begin{enumerate}
		\item 
		Comme $a \geqslant 0$ et $x^2 \geqslant 0$, on a $ax^2 \geqslant 0$, donc $ax^2 + 1 \geqslant 1 > 0$ : la fonction $f$ est donc dérivable sur $\R$ et sur cet intervalle $f'(x) = \dfrac{2ax}{ax^2 + 1}$.
		\item 
		Les résultats du 1. peuvent s'écrire :
		
$\left\{\begin{array}{l c l}
f(1)&=&3\\
f'(1)&=&1
\end{array}\right. \iff \left\{\begin{array}{l c l}
\ln (a + 1) + b&=&3\\
\dfrac{2a}{a + 1}&=&1
\end{array}\right.$.

La deuxième équation donne $2a = a + 1 \iff a = 1$ et en reportant dans la première :

$\ln (1 + 1) + b = 3 \iff b = 3 - \ln 2$.

On a donc sur $\R, \: f(x) = \ln \left(x^2 + 1 \right) + 3 - \ln 2$.

	\end{enumerate}
\end{enumerate} 

\bigskip

\textbf{Partie B}

\medskip

On admet que la fonction $f$ est définie sur $\R$ par

\[f(x) = \ln \left(x^2 + 1\right) + 3 - \ln (2).\]

\medskip

\begin{enumerate}
\item 
Quel que soit le réel $x, \: f(- x) = \ln \left[(- x)^2 + 1\right] + 3 - \ln 2 = \ln \left(x^2 + 1\right) + 3 - \ln (2) = f(x)$. La fonction $f$ est donc paire (la représentation graphique de $f$ est donc symétrique autour de l'axe des ordonnées).
\item 
On a $\displaystyle\lim_{x \to + \infty} x^2 = + \infty$ d'où $\displaystyle\lim_{x \to + \infty} x^2 + 1  = + \infty$ et par composition $\displaystyle\lim_{x \to + \infty} \ln \left(ax^2 + 1 \right) = + \infty$ et enfin $\displaystyle\lim_{x \to + \infty} f(x) = + \infty$.

La fonction étant paire $\displaystyle\lim_{x \to - \infty} f(x) = \displaystyle\lim_{x \to + \infty} f(x) = + \infty$.
\item 
Comme $x^2 + 1 > 0$ quel que soit le réel $x$, la fonction $f$ est dérivable sur $\R$ et sur cet intervalle : $f'(x) = \dfrac{2x}{x^2 + 1}.$

Le dénominateur étant supérieur à zéro le signe de $f('x)$ est donc celui de $2x$, donc :

$f'(x) < 0$ sur $\R_{-}^*$ et $f'(x) > 0$ sur $\R_{+}^*$.
Conclusion $f$ est décroissante sur $\R_{-}^*$ et croissante sur $\R_{+}^*$.

Le nombre $f(0) = \ln 1 + 3 - \ln 2 = 3 - \ln 2$ est donc le minimum de la fonction sur $\R$. D'où le tableau de variations :

\begin{center}
\psset{unit=1cm}
\begin{pspicture}(7,3)
\psframe(7,3)\psline(0,2)(7,2)\psline(0,2.5)(7,2.5)\psline(1,0)(1,3)
\uput[u](0.5,2.4){$x$} \uput[u](1.5,2.4){$- \infty$} \uput[u](4,2.4){$0$} \uput[u](6.5,2.4){$+\infty$} 
\uput[u](0.5,1.9){$f'(x)$}\uput[u](2.5,1.9){$-$}\uput[u](4,1.9){$0$}\uput[u](5.5,1.9){$+$}
\uput[d](1.5,2){$+ \infty$}\uput[d](6.5,2){$+ \infty$}\uput[u](4,0){$3 - \ln 2$}
\rput(0.5,1){$f(x)$}
\psline{->}(2,1.5)(3.5,0.5) \psline{->}(4.5,0.5)(6.5,1.5)
\end{pspicture}
\end{center}




\item 
D'après le tableau de variations l'équation $f(x) = k$ admet deux solutions si $k > 3 - \ln 2$.
\item 
$f(x) = 3 + \ln 2 \iff \ln \left(x^2 + 1\right) + 3 - \ln (2) = 3 + \ln (2) \iff \ln \left(x^2 + 1\right) = 2\ln (2) \iff \ln \left(x^2 + 1\right) = \ln 4 \iff x^2 + 1 = 4$ (par croissance de la fonction logarithme), soit $x^2 = 3$, d'où deux solutions $S = \left\{- \sqrt{3}~;~ \sqrt{3}\right\}$.
\end{enumerate}

\bigskip

\textbf{Partie C}

\medskip

On rappelle que la fonction $f$ est définie sur $R$ par $f(x) = \ln \left(x^2 + 1\right) + 3 - \ln (2)$.

\medskip

\begin{enumerate}
\item 

Il semble qu'il y ait deux points d'inflexion aux points d'abscisses $- 1$ et 1.
\item 
Comme $f'(x) = \dfrac{2x}{x^2 + 1}$ soit le quotient de deux fonctions dérivables sur $\R$, le dénominateur étant non nul ; $f'$ est donc dérivable sur $\R$ et :

$f''(x) = \dfrac{2\left(x^2 + 1\right) - 2x \times 2x}{\left(x^2 + 1\right)^2} = \dfrac{2x^2 + 2 - 4x^2}{\left(x^2 + 1\right)^2} = \dfrac{2\left(1 - x^2\right)}{\left(x^2 + 1\right)^2}$.
\item 
On a donc $f''(x) = 0 \iff 1 - x^2 = 0 \iff \left\{\begin{array}{l c  l}
1 +x&=&0\\
1 - x&=&
\end{array}\right.\iff \left\{\begin{array}{l c l}
x&=&- 1\\
x&=&1
\end{array}\right.$ Donc $S = \{-1~;~1\}$.

La dérivée seconde est positive quand le trinôme $1 - x^2$ est positif soit sur l'intervalle $]-1~;~1[$. Donc la fonction $f$ est convexe sur $]-1~;~1[$.
\end{enumerate}

\end{document}
